\section{为什么要研究 Agent Framework}

\subsection{从单 Agent 到多 Agent 协作}

\begin{figure}[H]
\centering
\includegraphics[page=6,width=0.9\textwidth]{slides.pdf}
\caption{导论把多 Agent 协作描述为复杂任务中的``分工机制''，不同模块负责不同子任务。}
\end{figure}

当任务跨度足够长、子任务差异足够大时，单个 Agent 往往会在上下文拥挤、职责混杂和推理漂移中失效。课程因此把多 Agent 协作列为核心主题之一：
\begin{itemize}
    \item 专门的 planner 负责任务拆解；
    \item specialized executor 负责代码、检索、表格分析等局部操作；
    \item verifier / reviewer 负责检查事实性、一致性和安全性。
\end{itemize}

\begin{importantbox}{Framework 的真正价值}
Agent framework 的价值不在于多套 API，而在于它为复杂任务提供了稳定的组织方式：如何拆分模块、如何传递上下文、如何让多个角色协作、如何把失败恢复写进执行流程。
\end{importantbox}

\subsection{应用机会与评测压力}

\begin{figure}[H]
\centering
\includegraphics[page=8,width=0.9\textwidth]{slides.pdf}
\caption{课程把代码生成、工作流自动化、个人助手和机器人并列展示，说明 Agent 的应用边界正在快速扩张。}
\end{figure}

\begin{figure}[H]
\centering
\includegraphics[page=9,width=0.9\textwidth]{slides.pdf}
\caption{GAIA、WebArena、SWE-Bench 等基准构成了 Agent 研究的早期评测坐标。}
\end{figure}

课程一开始就把应用和 benchmark 放在一起，是因为 Agent 系统很容易做出 demo，却很难稳定上线。导论中列出的挑战包括：
\begin{itemize}
    \item 复杂任务中的推理和规划不稳；
    \item 长时程任务中对环境反馈的利用效率不高；
    \item 多模态理解和 world model 仍不成熟；
    \item 多 Agent 协作、理论解释和安全隐私尚未形成稳固框架。
\end{itemize}

\begin{warningbox}{导论里最重要的风险提醒}
Agent 系统的失败通常不是``答错一道题''，而是长期任务中的累积偏差：错误计划、错误调用工具、错误读写外部环境，最后把小问题放大成系统性故障。也因此，部署 Agent 比部署普通 LLM 应用更强调可观测性和评测。
\end{warningbox}

\subsection{本章小结}

研究 Agent framework 的原因很实际：它既决定系统能否处理真实任务，也决定这些系统能否被稳定评估、调试和部署。

